\documentclass[10pt,a4paper]{amsart}
\usepackage[margin=19mm]{geometry}
\usepackage{amsmath,amssymb,mathtools}
\usepackage[scr=rsfs,cal=euler]{mathalfa}
\usepackage{xfrac}
\usepackage[unicode,hidelinks,bookmarks=false]{hyperref}
\newcommand{\ie}{i.e.\ }
\newcommand{\cf}{cf.\ }
\newcommand{\resp}{respectively\ }
\newcommand{\Subsection}{\subsection}
\providecommand{\nobreakdash}{\nobreak\hbox{-}}
\setlength{\emergencystretch}{2em}
\multlinegap0pt
\usepackage{bm}
\usepackage{bbm}
\usepackage{dsfont}
\usepackage{xspace}
\usepackage{cancel}
\usepackage{mathtools}
\usepackage{amsthm}
\makeatletter
\@ifundefined{theorem}{\newtheorem{theorem}{Theorem}}{}
\@ifundefined{assumption}{\newtheorem{assumption}{Assumption}}{}
\@ifundefined{lemma}{\newtheorem{lemma}[theorem]{Lemma}}{}
\@ifundefined{thm}{\newtheorem{thm}{Theorem}}{}
\makeatother
\title[Convergence of two-timescale gradient descent ascent dynamic]{Convergence of two-timescale gradient descent ascent dynamics: finite-dimensional and mean-field perspectives}
\author{Jing An, Jianfeng Lu}
\date{Source: arXiv:2501.17122v3 (CC BY 4.0)}
\begin{document}
\maketitle

\noindent\emph{Source and licence.} Convergence of two-timescale gradient descent ascent dynamics: finite-dimensional and mean-field perspectives. Version arXiv:2501.17122v3 (CC BY 4.0), distributed under the Creative Commons Attribution 4.0 International licence, as recorded in the arXiv licence metadata for this exact version and shown on the abstract page https://arxiv.org/abs/2501.17122v3. Full rights notice: licenses/arxiv-2501.17122-CC-BY-4.0.txt.\par\smallskip

\noindent\textit{Editorial scope.} Counted unit: the Thm whose source label is \texttt{thm:hypocoercivity} in the article (source0.tex lines 546--551, source label \texttt{thm:hypocoercivity}), delivered together with the article's own complete original proof of it (source0.tex lines 552--709, one \texttt{proof} environment of 158 lines). No mathematics is written, repaired or re-derived: statement and proof are the article's own text line for line. Required context retained uncounted: eqn:time\_s (lines 412--414); operator:M (lines 436--438); assump:ag (lines 440--444); assump:D (lines 460--474); assump:L (lines 476--484); lem:bounds (lines 487--492); assump:MAupper (lines 519--539); eqn:normequiv (lines 541--543). Every definition, display and label of those lines is retained unchanged. Omitted from this package and not redistributed: 1--411, 415--435, 439--439, 445--459, 475--475, 485--486, 493--518, 540--540, 544--545, 710--1393. Nothing inside a retained excerpt is omitted or altered: every retained body is delivered exactly as the article's lines are written, so no omission receipt of any kind accompanies this package. Citations: the retained excerpts carry no citation command at all, so no bibliography entry of the article is transcribed and no reference is redistributed. Reference adaptation: none was needed and none was made; every reference inside the delivered unit resolves inside it, so no unresolved reference or citation is produced. Printed slips kept literal: no printed word, symbol, hypothesis or number of the counted statement or of its proof was altered. Layout and class: the article's own class is \texttt{article} with packages including color, xcolor, graphicx, mathrsfs, amssymb, bbm, amsthm, amsmath, geometry, caption, numbysec, subcaption, float, mdframed; the wrapper is the lane's pinned \texttt{amsart} editorial wrapper at 10pt on A4, which restates the theorem-like environments the retained text needs and adds the article's own macro definitions that the retained text needs (none), with extra wrapper packages (bm, bbm, dsfont, xspace, cancel, mathtools). The delivered numbering is the unit's own. Assets: no retained span contains a figure environment, a graphics-inclusion command, a PDF-page inclusion, a table or a TikZ picture; no image file is copied into this package, the package declares no asset dependency and no image file is redistributed with it. Licence: version arXiv:2501.17122v3 of this article is distributed under the Creative Commons Attribution 4.0 International licence, as recorded in the arXiv licence metadata for this exact version and shown on the abstract page https://arxiv.org/abs/2501.17122v3. A full rights notice is delivered at licenses/arxiv-2501.17122-CC-BY-4.0.txt. Characters: every retained excerpt is pure ASCII, so the delivered engine, XeLaTeX with its default Latin Modern text font, reports no missing character.
\begin{equation}\label{eqn:time_s}
    \frac{d}{d s} \phi = -\frac{1}{\sqrt{\eta}}D \phi- L\phi.
\end{equation}

\begin{equation}\label{operator:M}
    M :=- ( I+(L\Pi)^{\top} L\Pi)^{-1}(L\Pi)^{\top}.
\end{equation}

\begin{assumption}\label{assump:ag} The projection operator $\Pi$ satisfies
   \begin{equation}
    \Pi L \Pi = 0.
\end{equation} 
\end{assumption}

\begin{assumption}[microscopic coercivity]\label{assump:D}
 Given $S = \Bigl[\begin{smallmatrix}
            Q & 0\\
            0 &  R
        \end{smallmatrix}\Bigr]$, let $\lambda_Q$ and $\lambda_R$ be the least eigenvalues of $Q$ and $R$ respectively.  If $\Pi=\Pi_Q$, we have
        
        \begin{equation}
            \langle S\phi, \phi\rangle\geq \lambda_Q \|(I-\Pi_Q)\phi\|^2,
        \end{equation}
 and if $\Pi=\Pi_R$, we have
        
        \begin{equation}
           \langle S\phi, \phi\rangle\geq \lambda_R \|(I-\Pi_R)\phi\|^2.
        \end{equation}
\end{assumption}

\begin{assumption}[macroscopic coercivity]\label{assump:L}
Let $\lambda_L$ be the minimum eigenvalue of $L^{\top}L = \Bigl[ \begin{smallmatrix}
            PP^\top & 0\\
            0 & P^\top P
        \end{smallmatrix}\Bigr]$, then
        \begin{equation}
            \|L\Pi\phi\|^2\geq \lambda_L \|\Pi\phi\|^2.
        \end{equation}
\end{assumption}

\begin{lemma}\label{lem:bounds}
    Given the operator $M$ in (\ref{operator:M}) and Assumption \ref{assump:ag}, we have the bounds
    \begin{equation}
         \|M\phi\|\leq \frac{1}{2}\|(I-\Pi)\phi\|,\quad \|LM\phi\|\leq\|(I-\Pi)\phi\|.
    \end{equation}
\end{lemma}

\begin{assumption}\label{assump:MAupper}
%     Moreover, we assume that
% \begin{equation}
%     \langle MA^{\top}A\phi, \phi\rangle\leq \Big(\Lambda_Q\|(I-\Pi_Q)\phi\|+\eta \Lambda_R\|(I-\Pi_R)\Pi_Q\phi\|\Big)\|\phi\|.
% \end{equation}
 Given $S = \Bigl[\begin{smallmatrix}
            Q & 0\\
            0 &  R
        \end{smallmatrix}\Bigr]$, if $\Pi=\Pi_Q$, there exists $\Lambda_Q>0$ such that      
        \begin{equation}
           \langle M S\phi, \phi\rangle\leq \Lambda_Q\|(I-\Pi_Q)\phi\|\|\phi\|
        \end{equation}
 and if $\Pi=\Pi_R$, then there exists $\Lambda_R>0$ such that
        \begin{equation}
          \langle M S\phi, \phi\rangle\leq  \Lambda_R\|(I-\Pi_R)\phi\|\|\phi\|.
        \end{equation}
        Moreover, we can find a constant $C_M>0$ such that 
     \begin{equation}
              \langle M L(I-\Pi)\phi, \phi\rangle\leq C_M \|(I-\Pi)\phi\|\|\phi\|.
     \end{equation}
\end{assumption}

\begin{equation}\label{eqn:normequiv}
   \frac{1-\epsilon}{2}\|\phi\|^2\leq  H(\phi)\leq \frac{1+\epsilon}{2}\|\phi\|^2.
\end{equation}

\begin{thm}\label{thm:hypocoercivity}
With Assumptions \ref{assump:ag}, \ref{assump:D}, \ref{assump:L}, and \ref{assump:MAupper}, we can find explicit constants $\Lambda, C>0$, depending on $\epsilon$,  to get that for all $s \geq 0$,
\begin{equation}
    \|\phi(s)\|^2 \leq C \exp\Big(-\Lambda \min\Big\{\sqrt{\eta}, \frac{1}{\sqrt{\eta}}\Big\} s \Big)\|\phi_0\|^2.
\end{equation}
\end{thm}

\begin{proof}
\textbf{Case 1: $\eta\ll 1$.} In this case, we consider 
$\tilde D  := \bigl[\begin{smallmatrix}
            Q & 0\\
            0 & 0
        \end{smallmatrix}\bigr]$
as an approximation to the diffusion matrix $D= \bigl[\begin{smallmatrix}
            Q & 0\\
            0 & \eta R
        \end{smallmatrix}\bigr]$,
       We can find a constant $C_R>0$ such that
        \begin{equation}\label{eqn:err1}
            \|\tilde D- D\|_F\leq C_R \eta.
        \end{equation}
With $\tilde D$ in mind, we choose $\Pi = \Pi_Q$ satisfying Assumption \ref{assump:ag} that projects onto kernel of  $\tilde D$, and write $H(\phi) = H_Q(\phi)$ accordingly. Taking the derivative of $H(\phi)$ in time $s$, by the dynamics (\ref{eqn:time_s}), we have
\begin{equation}
    \begin{aligned}
        \frac{d}{ds} H_Q(\phi) &= \langle \phi, \phi_s\rangle-\epsilon\langle M\phi_s, \phi\rangle-\epsilon\langle M\phi, \phi_s\rangle\\
        &= -\frac{1}{\sqrt{\eta}}\langle \phi, D\phi\rangle + \frac{\epsilon}{\sqrt{\eta}}\langle MD\phi, \phi\rangle+\epsilon\langle M L\phi, \phi\rangle\\
        &\quad +\frac{\epsilon}{\sqrt{\eta}}\langle M\phi, D\phi\rangle+\epsilon\langle M\phi, L \phi\rangle.
    \end{aligned}
\end{equation}
For $\langle \phi, D\phi\rangle$, we utilize Assumption \ref{assump:D} and (\ref{eqn:err1}) to obtain that
\begin{equation}
\begin{aligned}
    -\frac{1}{\sqrt{\eta}}\langle   D \phi, \phi\rangle& \leq  -\frac{1}{\sqrt{\eta}}\langle  \tilde D \phi, \phi\rangle +\frac{1}{\sqrt{\eta}}\big|\langle (D - \tilde D)\phi, \phi\rangle \big|\\
    &\leq -\frac{\lambda_Q}{\sqrt{\eta}}\|(I-\Pi_Q)\phi\|^2+C_R \sqrt{\eta} \|\phi\|^2.
\end{aligned}
\end{equation}
For $\langle MD\phi, \phi\rangle$, by the property that $\Pi_Q M= M$ and Assumption \ref{assump:ag}, we observe that
\begin{equation}
   \langle M \tilde{D}\phi, \phi \rangle = \langle M\phi, \tilde{D}\phi\rangle = \langle \Pi_Q M\phi, \tilde{D}\phi\rangle=\langle \tilde{D} \Pi_Q M\phi, \phi\rangle = 0.
\end{equation}
Thus
\begin{equation}
    \frac{\epsilon}{\sqrt{\eta}}\langle MD\phi, \phi\rangle\leq \frac{\epsilon}{\sqrt{\eta}}\big|\langle M(D-\tilde{D})\phi, \phi\rangle\big|\leq \epsilon C_R\sqrt{\eta}\|\phi\|^2.
\end{equation}
As for $\langle M L\phi, \phi\rangle$, we may split
\begin{equation}
    \begin{aligned}
        \langle M L\phi, \phi\rangle = \langle M L\Pi\phi, \phi\rangle+ \langle M L(I-\Pi)\phi, \phi\rangle,
    \end{aligned}
\end{equation}
and treat the first part by using Assumption \ref{assump:L}:
\begin{equation}
    \langle M L\Pi\phi, \phi\rangle = -\langle( I+(L\Pi)^{\top} L\Pi)^{-1}(L\Pi)^{\top}(L\Pi)\phi, \phi\rangle \leq -\frac{\lambda_L}{1+\lambda_L} \|\Pi\phi\|^2 = -\frac{\lambda_L}{1+\lambda_L} \|\Pi_Q\phi\|^2.
\end{equation}
As for the second part, we have the bound  by  Assumption \ref{assump:MAupper},
\begin{equation}
     \langle M L(I-\Pi_Q)\phi, \phi\rangle\leq C_M \|(I-\Pi_Q)\phi\|\|\phi\|.
\end{equation}
For $\langle M\phi, D\phi\rangle$, we then use bounds in Lemma \ref{lem:bounds} and (\ref{eqn:err1}) to get that
\begin{equation}
\begin{aligned}
   \frac{\epsilon}{\sqrt{\eta}} \langle M\phi, D\phi\rangle&=\frac{\epsilon}{\sqrt{\eta}} \langle M\phi, \tilde{D}\phi\rangle+\frac{\epsilon}{\sqrt{\eta}} \langle M\phi, (D-\tilde{D})\phi\rangle\\
   &\leq \frac{\epsilon}{\sqrt{\eta}}\|M\phi\|\Big(\|\tilde{D}\phi\|+\|(D-\tilde{D})\phi\|\Big) \\
   &\leq \frac{\epsilon\Lambda_Q}{2\sqrt{\eta}}\|(I-\Pi_Q)\phi\|\|\phi\| + \frac{\epsilon C_R\sqrt{\eta}}{2} \|(I-\Pi_Q)\phi\| \|\phi\|.
\end{aligned}
\end{equation}
Lastly,  for $\langle M\phi, L \phi\rangle$, by Lemma \ref{lem:bounds}, we obtain that
\begin{equation}
    \langle M\phi, L \phi\rangle = -\langle LM\phi,  \phi\rangle\leq  \|(I-\Pi_Q)\phi\|\|\phi\|.
\end{equation}
Combining all bounds above together, we have that
\begin{equation}
\begin{aligned}
    \frac{d}{ds} H_Q(\phi) &\leq -\frac{\lambda_Q}{\sqrt{\eta}}\|(I-\Pi_Q)\phi\|^2-\frac{\epsilon\lambda_L}{1+\lambda_L} \|\Pi_Q\phi\|^2\\
    &\quad + \epsilon\Big(\frac{\Lambda_Q}{2\sqrt{\eta}}+\frac{C_R\sqrt{\eta}}{2}+1+C_M\Big)\|(I-\Pi_Q)\phi\|\|\phi\| + (1+\epsilon) C_R \sqrt{\eta} \|\phi\|^2\\
    & = -\begin{bmatrix}
            \|(I-\Pi_Q)\phi\| \\
            \|\Pi_Q\phi\|
        \end{bmatrix}^\top \begin{bmatrix}
            S_{++} & S_{+-}/2\\
          S_{+-}/2 & S_{--}
        \end{bmatrix} \begin{bmatrix}
            \|(I-\Pi_Q)\phi\| \\
            \|\Pi_Q\phi\| 
        \end{bmatrix}+(1+\epsilon)C_R \sqrt{\eta}\|\phi\|^2,
\end{aligned}
\end{equation}
with 
\begin{equation}
    S_{++} = \frac{\lambda_Q}{\sqrt{\eta}}, \quad S_{--} = \frac{\epsilon\lambda_L}{1+\lambda_L} , \quad S_{+-} = - \epsilon\Big(\frac{\Lambda_Q}{2\sqrt{\eta}}+\frac{C_R\sqrt{\eta}}{2}+1+C_M\Big)
\end{equation}
and the smallest eigenvalue of the matrix $\bigl[\begin{smallmatrix}
            S_{++} & S_{+-}/2\\
          S_{+-}/2 & S_{--}
        \end{smallmatrix} \bigr]$ is
\begin{equation}
    \lambda_{\min} = \frac{1}{2}\Big(S_{++} + S_{--}-\sqrt{(S_{++}-S_{--})^2+(S_{+-})^2}\Big).
\end{equation}
Note that we consider the regime where $\eta\ll1$. In order to make $\lambda_{\min}>0$,  we take $\epsilon \sim \sqrt{\eta}$ so that for sufficiently small $\eta$, we can find a constant $c>0$ to get
\begin{equation}
     \frac{d}{ds} H_Q(\phi) \leq -c\sqrt{\eta} \|\phi\|^2\leq -\frac{2c\sqrt{\eta}}{1+\epsilon} H_Q(\phi)
\end{equation}
by the norm equivalence (\ref{eqn:normequiv}), which also implies that for $\eta\ll 1$,
\begin{equation}
    H_Q(\phi(s))\leq \exp\Big(\frac{-2c\sqrt{\eta}s}{1+\epsilon}\Big) H_Q(\phi_0) \quad \Longrightarrow \quad \|\phi(s)\|^2 \leq \frac{1+\epsilon}{1-\epsilon}\exp\Big(\frac{-2c\sqrt{\eta}s}{1+\epsilon}\Big) \|\phi_0\|^2.
\end{equation}
\textbf{Case 2: $\eta\gg 1$.}
We first rewrite the dynamics (\ref{eqn:time_s}) as
\begin{equation}
     \phi_s = -\sqrt{\eta }B \phi- L\phi
\end{equation}
with the re-scaled diffusion matrix  
$B= \frac{1}{\eta} D=  \bigl[ \begin{smallmatrix}
            Q/\eta & 0\\
            0 & R
 \end{smallmatrix}\bigr]$. We take 
$\tilde B  := \bigl[\begin{smallmatrix}
            0 & 0\\
            0 & R
        \end{smallmatrix}\bigr]$ to compare with $B$, and one can find a constant $C_Q>0$ such that 
        \begin{equation}
            \|\tilde B- B\|_F\leq C_Q \eta^{-1}.
        \end{equation}
With $\tilde B$ in mind, we choose $\Pi = \Pi_R$ satisfying Assumption \ref{assump:ag} that projects onto kernel of  $\tilde B$, and write $H(\phi) = H_R(\phi)$ accordingly. Taking derivative of $H(\phi)$ in time $s$, we then get
\begin{equation}
    \begin{aligned}
        \frac{d}{ds} H_R(\phi) &= \langle \phi, \phi_s\rangle-\epsilon\langle M\phi_s, \phi\rangle-\epsilon\langle M\phi, \phi_s\rangle\\
        &= -\sqrt{\eta}\langle \phi, B\phi\rangle + \epsilon\sqrt{\eta}\langle M B\phi, \phi\rangle+\epsilon\langle M L\phi, \phi\rangle\\
        &\quad +\epsilon\sqrt{\eta}\langle M\phi, B\phi\rangle+\epsilon\langle M\phi, L \phi\rangle.
    \end{aligned}
\end{equation}
The estimates of each term are similar to what we have done in Case 1. Omitting those details, we will eventually get the bound that
\begin{equation}
\begin{aligned}
    \frac{d}{ds} H_R(\phi) &\leq -\lambda_R\sqrt{\eta}\|(I-\Pi_R)\phi\|^2-\frac{\epsilon\lambda_L}{1+\lambda_L} \|\Pi_R\phi\|^2\\
    &\quad+\epsilon\Big(\frac{\Lambda_R \sqrt{\eta}}{2}+\frac{C_R}{2\sqrt{\eta}}+1+C_M\Big)\|(I-\Pi_R)\phi\|\|\phi\|
  + \frac{(1+\epsilon)C_Q}{\sqrt{\eta}} \|\phi\|^2\\
    & = -\begin{bmatrix}
            \|(I-\Pi_Q)\phi\| \\
            \|\Pi_Q\phi\|
        \end{bmatrix}^\top \begin{bmatrix}
            S_{++} & S_{+-}/2\\
          S_{+-}/2 & S_{--}
        \end{bmatrix} \begin{bmatrix}
            \|(I-\Pi_Q)\phi\| \\
            \|\Pi_Q\phi\| 
        \end{bmatrix}+ \frac{(1+\epsilon) C_Q}{\sqrt{\eta}} \|\phi\|^2,
\end{aligned}
\end{equation}
with \begin{equation}
    S_{++} = \lambda_R\sqrt{\eta}, \quad S_{--} = \frac{\epsilon\lambda_L}{1+\lambda_L} , \quad S_{+-} = - \epsilon\Big(\frac{\Lambda_R \sqrt{\eta}}{2}+\frac{C_R}{2\sqrt{\eta}}+1+C_M\Big)
\end{equation}
and the smallest eigenvalue of this matrix is
\begin{equation}
    \lambda_{\min} = \frac{1}{2}\Big(S_{++} + S_{--}-\sqrt{(S_{++}-S_{--})^2+(S_{+-})^2}\Big).
\end{equation}
Since we consider the regime where $\eta\gg 1$, in order to make $\lambda_{\min}>0$ we take $\epsilon \sim 1/\sqrt{\eta}$ so that there exists a constant $c'>0$,
\begin{equation}
     \frac{d}{ds} H_R(\phi) \leq -\frac{c'}{\sqrt{\eta}} \|\phi\|^2\leq -\frac{2c'}{(1+\epsilon)\sqrt{\eta}} H_R(\phi).
\end{equation}
Accompanied with the norm equivalence (\ref{eqn:normequiv}), we conclude that for $\eta\gg1$, 
\begin{equation}
    H_R(\phi(s))\leq \exp\Big(\frac{-2c's}{(1+\epsilon)\sqrt{\eta}}\Big) H_R(\phi_0) \quad \Longrightarrow \quad \|\phi(s)\|^2 \leq \frac{1+\epsilon}{1-\epsilon}\exp\Big(\frac{-2c's}{(1+\epsilon)\sqrt{\eta}}\Big) \|\phi_0\|^2.
\end{equation}
\end{proof}

\end{document}
